% Paste this environment into the IEEEtran source.
% The preamble must already \input the sibling iscafig.tex.
\begin{figure}[t]
\centering
\begin{tikzpicture}[iscafig]
  \node[lbl, anchor=west] at (0,1.05) {one register, 36 bits};
  \node[fldacc, anchor=west, minimum width=1.15cm, minimum height=7.2mm] (tag) at (0,0)
    {tag};
  \node[fldcore, anchor=west, minimum width=6.55cm, minimum height=7.2mm] (pay) at (1.15,0)
    {payload};
  \node[idx, anchor=south] at (0.55,0.42) {35\,:\,32};
  \node[idx, anchor=south] at (4.4,0.42) {31\,:\,0};

  \node[lbl, anchor=west] at (0,-0.85) {values that occupy more than one register};
  \node[fldcore, minimum width=3.55cm, minimum height=6.6mm, anchor=west] (i0) at (0,-1.35)
    {\texttt{INT64} low 32};
  \node[fldacc, minimum width=3.55cm, minimum height=6.6mm, anchor=west] (i1) at (3.7,-1.35)
    {\texttt{CONT} high 32};
  \node[fldcore, minimum width=3.55cm, minimum height=6.6mm, anchor=west] (c0) at (0,-2.15)
    {\texttt{COMPLEX} real};
  \node[fldacc, minimum width=3.55cm, minimum height=6.6mm, anchor=west] (c1) at (3.7,-2.15)
    {\texttt{CONT} imag.\ float};
  \node[fldhot, minimum width=1.75cm, minimum height=6.6mm, anchor=west] (s0) at (0,-2.95)
    {\scriptsize\texttt{SHORT\_STR}};
  \node[fldacc, minimum width=1.75cm, minimum height=6.6mm, anchor=west] (s1) at (1.85,-2.95)
    {\scriptsize\texttt{CONT}};
  \node[fldacc, minimum width=1.75cm, minimum height=6.6mm, anchor=west] (s2) at (3.70,-2.95)
    {\scriptsize\texttt{CONT}};
  \node[fldacc, minimum width=1.75cm, minimum height=6.6mm, anchor=west] (s3) at (5.55,-2.95)
    {\scriptsize\texttt{CONT}};
  \node[lbl, anchor=west, align=left] at (0,-3.7)
    {\texttt{SHORT\_STR} always uses four registers, at most 15 bytes.\\
     \texttt{CONT} is not a value on its own.};
\end{tikzpicture}
\caption{Tagged register.
Every other value in the tag table occupies one register.
A \texttt{MUT\_COLLEC} payload keeps its kind in bits\,[1:0] and the address is those bits cleared.
Pointer tags are \texttt{STR}, \texttt{TUPLE}, \texttt{MUT\_COLLEC}, \texttt{RANGE}, \texttt{ITER}, \texttt{BYTES}, and \texttt{OBJECT}.}
\label{fig:value}
\end{figure}
